\documentclass[11pt]{article}
\usepackage[margin=1in]{geometry}
\usepackage{amsmath,amssymb,amsthm}
\usepackage[T1]{fontenc}
\usepackage{lmodern}
\usepackage{xcolor}
\usepackage{booktabs}
\usepackage{array}
\usepackage[hidelinks]{hyperref}
\usepackage{url}

\newtheorem{theorem}{Theorem}
\newtheorem{lemma}[theorem]{Lemma}
\theoremstyle{definition}
\newtheorem{definition}[theorem]{Definition}
\theoremstyle{remark}
\newtheorem{remark}[theorem]{Remark}

\newcommand{\lean}[1]{\texttt{#1}}
% A declaration of one of our developments; every \decl{...} must be in that repository's gate.sh
% REQUIRED list (checked by the name gate).
\DeclareUrlCommand\decl{\urlstyle{tt}\def\UrlBreaks{\do\_\do\.}}
% A Mathlib declaration (pinned revision 0df444a3); not one of ours.
\DeclareUrlCommand\mdecl{\urlstyle{tt}\def\UrlBreaks{\do\_\do\.}}
\newcommand{\RR}{\mathbb{R}}
\newcommand{\CC}{\mathbb{C}}
\newcommand{\PH}[1]{\textcolor{red}{\textbf{[PLACEHOLDER:} #1\textbf{]}}}

% ---------------------------------------------------------------------------------------------
% PINS. Filled in when the repositories are published; until then \src prints file:line only.
% ---------------------------------------------------------------------------------------------
\newcommand{\commitof}[1]{\csname commitof:#1\endcsname}
\expandafter\def\csname commitof:quadnewton\endcsname{48e94cf944b3dfbcca7358360a064e39b6336d98}
\expandafter\def\csname commitof:quadmiquel\endcsname{011c9ecd5276b7abaf2b42a8f5ba463196df123d}
\newcommand{\repo}[1]{https://github.com/jpbald93/#1}
\newcommand{\src}[4]{\ifcsname commitof:#1\endcsname
  \href{\repo{#1}/blob/\commitof{#1}/#2/#3.lean\#L#4}{\texttt{#3}:#4}\else\texttt{#3}:#4\fi}

\title{The complete quadrilateral in Lean~4:\\
Newton--Gauss, Bodenmiller, the Miquel point, the Steiner circle and the Simson line}
\author{Josh Bald\\ \small Independent researcher, Ontario, Canada\\ \small\texttt{jpbald93@gmail.com}\\
\small ORCID: \texttt{0009-0002-1317-6489}
\thanks{\copyright~2026 Josh Bald. Licensed under CC BY 4.0 (\url{https://creativecommons.org/licenses/by/4.0/}). The accompanying Lean code is licensed under the Apache License 2.0.}}
\date{15 October 2026}

\emergencystretch=3em
\begin{document}
\sloppy
\maketitle

\begin{abstract}
Four lines in general position form a complete quadrilateral: six vertices, three diagonals, and
four triangles, one for each line omitted. We give machine-checked proofs, in Lean~4 with Mathlib,
of eight classical theorems about this configuration, all from the single hypothesis that the lines
are in general position. On the side of the diagonals: the midpoints of the three diagonals are
collinear (the Newton--Gauss line); the orthocentres of the four triangles are collinear; every
orthocentre has equal power with respect to the three circles on the diagonals as diameters (the
Gauss--Bodenmiller theorem); and the line of the orthocentres is perpendicular to the Newton--Gauss
line. On the side of the circles: the circumcircles of the four triangles pass through a common point
$M$, given by an explicit formula (the Miquel point); the four circumcentres and $M$ lie on one circle
(the Steiner, or Miquel, circle); the feet of the perpendiculars from $M$ to the four lines are
collinear (the Simson line of the quadrilateral); and that line is parallel to the line of the four
orthocentres, as Mathlib affine subspaces. The two halves are separate developments that share one file of definitions. We also
show that two of the four orthocentres can coincide in general position (two of them coincide exactly
when two of the lines are perpendicular), so the orthocentre line is determined only because the four
are never all equal, which we prove. The local definitions (line, intersection, circumcentre,
orthocentre, foot of a perpendicular, power of a point, general position) are proved equal to
Mathlib's. All proofs use only Lean's
three standard axioms. In the sources we searched, the Newton--Gauss line and the four-line Miquel
theorem already have formal proofs (in HOL~Light and Rocq); we found none of the other six in any
system, and none of the eight in Lean.
\end{abstract}

\section{Introduction}

A \emph{complete quadrilateral} is a set of four lines in the plane, no two parallel and no three
through a point. Its six intersection points are its \emph{vertices}; the three segments joining the
pairs of vertices that do not lie on a common line are its \emph{diagonals}; and omitting one line at a
time leaves four triangles. The configuration carries a remarkable family of classical theorems,
found between Gauss and Steiner in the first third of the nineteenth century
\cite{steiner1828,coolidge1916,johnson1929}. Two of them concern points: the midpoints of the diagonals
lie on a line, and so do the four orthocentres. Two concern circles: the circles on the diagonals as
diameters are coaxal, and the four circumcircles pass through one point. The rest tie the two
families together. We formalise eight of them (Table~\ref{tab:thms}).

\begin{table}[h]
\centering\small
\begin{tabular}{llll}
\toprule
& theorem & Lean & source (pinned)\\
\midrule
\ref{thm:newton} & midpoints of the diagonals collinear & \decl{QuadNewton.newton_gauss} & \src{quadnewton}{QuadNewton}{Quad}{53}\\
\ref{thm:ortho} & four orthocentres collinear (not all equal) & \decl{QuadNewton.orthocentres_collinear} & \src{quadnewton}{QuadNewton}{Quad}{79}\\
\ref{thm:boden} & diagonal circles coaxal (equal power) & \decl{QuadNewton.bodenmiller} & \src{quadnewton}{QuadNewton}{Quad}{135}\\
\ref{thm:perp} & orthocentre line $\perp$ Newton--Gauss line & \decl{QuadNewton.orthocentre_line_perp} & \src{quadnewton}{QuadNewton}{Quad}{155}\\
\ref{thm:miquel} & four circumcircles concur at $M$ & \decl{QuadMiquel.miquel} & \src{quadmiquel}{QuadMiquel}{Miquel}{162}\\
\ref{thm:steiner} & circumcentres and $M$ concyclic & \decl{QuadMiquel.steiner_circle} & \src{quadmiquel}{QuadMiquel}{Miquel}{200}\\
\ref{thm:simson} & feet from $M$ to the four lines collinear & \decl{QuadMiquel.simson} & \src{quadmiquel}{QuadMiquel}{Simson}{155}\\
\ref{thm:par} & that line $\parallel$ the orthocentre line & \decl{QuadMiquel.steiner_parallel_line} & \src{quadmiquel}{QuadMiquel}{Parallel}{241}\\
\bottomrule
\end{tabular}
\caption{The eight theorems. Each assumes only that the four lines are in general position.}\label{tab:thms}
\end{table}

\subsection{Contributions}
\begin{itemize}
\item One model of the complete quadrilateral in $\CC$, with explicit formulas for the intersection
point of two lines, the circumcentre, the orthocentre and the foot of a perpendicular
(Section~\ref{sec:model}), shared by two independent developments.
\item The eight theorems, each stated for arbitrary four lines with the single hypothesis
$\lean{GenPos}$ (Sections~\ref{sec:diag} and~\ref{sec:circ}). The collinearity and concyclicity
conclusions use Mathlib's \mdecl{Collinear} and \mdecl{EuclideanGeometry.Cospherical}.
\item An explicit formula for the Miquel point, as the centre of a spiral similarity
(Section~\ref{sec:circ}).
\item A small finding about the statements themselves: in general position two of the four
orthocentres can coincide; $H_3=H_4$ exactly when $L_1\perp L_2$ (\decl{QuadMiquel.ortho34_eq_iff}).
So ``the line of the orthocentres'' needs the fact, which we prove, that they are not all equal
(Section~\ref{sec:diag}), and Steiner's parallelism, stated for a fixed pair of orthocentres, needs
that pair to be distinct (Section~\ref{sec:circ}).
\item Bridges tying every local notion to Mathlib's (Section~\ref{sec:bridges}), and the uniqueness of
the Miquel point: any point on the circumcircles of $T_4$, $T_3$ and $T_1$ is $M$
(\decl{QuadMiquel.miquelPt_unique}); in particular the common point of all four is unique.
\end{itemize}

\paragraph{What is not claimed.} We prove each theorem for one fixed labelling of the four lines; the
statements are symmetric in content but we did not prove the symmetry. The diagonal circles are shown
to be coaxal in the sense that the four orthocentres have equal power with respect to all three, which
together with Theorem~\ref{thm:ortho} places the radical axis; we do not formalise coaxal pencils as
such.

\subsection{Availability}\label{sec:avail}
The proofs are in two Lean~4 developments on Lean v4.33.1 with Mathlib pinned at revision
\texttt{0df444a3}. File and line counts are of the library \lean{.lean} files, including the root
import file and counting every physical line; in total 14 files and 2{,}980 lines.
\begin{center}\small
\begin{tabular}{llrr}
\toprule
repository & content & files & lines\\
\midrule
\lean{quadnewton} & Theorems~\ref{thm:newton}--\ref{thm:perp} and bridges & 6 & 1{,}066\\
\lean{quadmiquel} & Theorems~\ref{thm:miquel}--\ref{thm:par} and bridges & 8 & 1{,}914\\
\bottomrule
\end{tabular}
\end{center}
The repositories are \url{https://github.com/jpbald93/quadnewton} at commit
\texttt{48e94cf944b3dfbcca7358360a064e39b6336d98} and \url{https://github.com/jpbald93/quadmiquel} at commit
\texttt{011c9ecd5276b7abaf2b42a8f5ba463196df123d}; Lean v4.33.1, Mathlib revision \texttt{0df444a360eaa60ab8c11dca51a86af692955474}
(pinned in \lean{lake-manifest.json}). Source links in this paper point to these commits.
Each repository has a gate script (\lean{gate.sh}) that rejects forbidden tokens, builds the library
and checks with \lean{\#print axioms} that every listed declaration depends only on
\lean{propext}, \lean{Classical.choice} and \lean{Quot.sound}, and a tampering suite
(\lean{tests/tamper.sh}) that plants defects the gate must reject.

\section{The model}\label{sec:model}

Both developments start from the same file \lean{Basic.lean}, identical up to the name of its
namespace. Points are complex numbers. For $a,b\in\CC$ let $\operatorname{cr}(a,b)=\operatorname{Im}(\bar a b)$,
the cross product of $a$ and $b$ as plane vectors. A line $L$ is a pair $(p,d)$ of a point and a
direction, with carrier $\{z:\operatorname{cr}(d,z-p)=0\}$. For lines $L=(p,d)$ and $M=(q,e)$,
\[
  \operatorname{meet}(L,M)=p+\frac{\operatorname{cr}(e,q-p)}{\operatorname{cr}(e,d)}\,d ,
\]
which is their intersection point when $\operatorname{cr}(e,d)\neq0$. Four lines $L_1,\dots,L_4$ are in
\emph{general position} ($\lean{GenPos}$) if their directions are nonzero, no two are parallel
($\operatorname{cr}(d_i,d_j)\ne0$), and no three pass through one point
($\operatorname{meet}(L_i,L_j)\notin L_k$ for the four triples). We write $A_{ij}=\operatorname{meet}(L_i,L_j)$
for the six vertices; the diagonals are $A_{12}A_{34}$, $A_{13}A_{24}$ and $A_{14}A_{23}$; and
$T_k$ is the triangle of the three vertices not on $L_k$: $T_4=(A_{12},A_{13},A_{23})$,
$T_3=(A_{12},A_{14},A_{24})$, $T_2=(A_{13},A_{14},A_{34})$, $T_1=(A_{23},A_{24},A_{34})$.

The circumcentre, the orthocentre and the foot of a perpendicular are given by explicit formulas:
\[
  \operatorname{circ}(a,b,c)=a+\frac{|b-a|^2(c-a)-|c-a|^2(b-a)}{\overline{(b-a)}(c-a)-(b-a)\overline{(c-a)}},\qquad
  \operatorname{ortho}(a,b,c)=a+b+c-2\operatorname{circ}(a,b,c),
\]
\[
  \operatorname{foot}(z,L)=p+\frac{\operatorname{Re}\big(\bar d(z-p)\big)}{|d|^2}\,d .
\]
The orthocentre is defined through Euler's relation $H=A+B+C-2O$. Section~\ref{sec:bridges} proves
that each of these is the corresponding Mathlib notion.

\paragraph{Normal coordinates.} Both developments first prove that general position puts the
configuration into a normal form. In \lean{quadnewton} (\decl{QuadNewton.normal_coords}): there are
$O,u,w\in\CC$ and real $s,t$ with $A_{12}=O$, $A_{13}=O+u$, $A_{23}=O+w$, $A_{14}=O+su$,
$A_{24}=O+tw$ and $A_{34}=O+\big(s(1-t)u+t(s-1)w\big)/(s-t)$, where $\operatorname{cr}(u,w)\neq0$,
$s,t\notin\{0,1\}$ and $s\neq t$. \lean{quadmiquel} proves the same normal form independently, in its own
notation. The central coordinate identities are then rational-function identities in
$u,w,\bar u,\bar w,s,t$ with those non-vanishing conditions, discharged principally by
\lean{field\_simp}, \lean{ring} and \lean{linear\_combination} with the conjugates treated as
independent variables; incidence and non-degeneracy lemmas connect them to the geometric conclusions. The two developments were written separately, by different agents on different
machines, from the shared definitions.

\section{The diagonals and the orthocentres}\label{sec:diag}

All four theorems of this section are in \lean{quadnewton}.

\begin{theorem}[Newton--Gauss line; \decl{QuadNewton.newton_gauss}]\label{thm:newton}
If $L_1,\dots,L_4$ are in general position, the midpoints of the three diagonals are collinear; and
the midpoints of $A_{12}A_{34}$ and $A_{13}A_{24}$ are distinct (\decl{QuadNewton.diag_midpoints_ne}).
\end{theorem}

\begin{theorem}[Orthocentre line; \decl{QuadNewton.orthocentres_collinear}, \decl{QuadNewton.orthocentres_not_all_eq}]\label{thm:ortho}
If $L_1,\dots,L_4$ are in general position, the orthocentres $H_1,\dots,H_4$ of $T_1,\dots,T_4$ are
collinear, and they are not all equal: $H_3\neq H_4$ or $H_2\neq H_4$.
\end{theorem}

The second part is needed, and cannot be strengthened to a fixed pair. For $L_1$ the real axis, $L_2$
the imaginary axis, $L_3$ the line through $1$ and $i$, and $L_4$ the line through $2$ and $-2i$, the
lines are in general position and $H_1=H_2=\tfrac32-\tfrac12i$, $H_3=H_4=0$ (checked in exact rational
arithmetic). So any one pair of orthocentres can coincide, and ``the line of the orthocentres'' is
determined only because the four are never all equal.

For points $P\neq Q$ let $\operatorname{pw}(z;P,Q)=|z-\tfrac{P+Q}2|^2-\tfrac14|P-Q|^2$, the power of $z$
with respect to the circle on $PQ$ as diameter (it is Mathlib's \mdecl{EuclideanGeometry.Sphere.power}
of that circle: \decl{QuadNewton.pw_eq_power}).

\begin{theorem}[Gauss--Bodenmiller; \decl{QuadNewton.bodenmiller}]\label{thm:boden}
If $L_1,\dots,L_4$ are in general position, each orthocentre $H_k$ has the same power with respect to
the three circles having the diagonals as diameters.
\end{theorem}

Since the centres of these circles (the diagonal midpoints) are pairwise distinct
(\decl{QuadNewton.bodenmiller_coaxal}, which also restates Theorem~\ref{thm:boden} with Mathlib's
power), the points of equal power with respect to two of them form a line, and Theorem~\ref{thm:boden} says all four orthocentres
lie on the radical axis of every pair; that is the classical statement that the three circles are
coaxal, with the orthocentre line as their radical axis.

\begin{theorem}[\decl{QuadNewton.orthocentre_line_perp}]\label{thm:perp}
If $L_1,\dots,L_4$ are in general position, then for any two orthocentres $H,H'$,
$\operatorname{Re}\big((m_2-m_1)\overline{(H'-H)}\big)=0$, where $m_1,m_2$ are the midpoints of
$A_{12}A_{34}$ and $A_{13}A_{24}$.
\end{theorem}
With $m_1\neq m_2$ (Theorem~\ref{thm:newton}) and two distinct orthocentres (Theorem~\ref{thm:ortho}),
this is the perpendicularity of the orthocentre line and the Newton--Gauss line. When $H=H'$ the
statement holds trivially. The real part is Mathlib's inner product on $\CC$
(\decl{QuadNewton.perp_iff_inner}; the inner-product form is \decl{QuadNewton.orthocentre_line_perp_inner}).

\section{The circles}\label{sec:circ}

All four theorems of this section are in \lean{quadmiquel}. The Miquel point is given explicitly:
\[
  M=A_{12}+\frac{(A_{13}-A_{12})(A_{24}-A_{12})-(A_{23}-A_{12})(A_{14}-A_{12})}{A_{13}+A_{24}-A_{23}-A_{14}}
\]
(\lean{QuadMiquel.miquelPt}), the centre of the spiral similarity taking $A_{13}\mapsto A_{14}$ and
$A_{23}\mapsto A_{24}$. General position makes the denominator nonzero. $M$ is the reflection of $A_{12}$ in the line
through the circumcentres of $T_4$ and $T_3$, so it is the second intersection point of those two
circles; it is none of the six vertices; and it is the only point on the circumcircles of $T_4$,
$T_3$ and $T_1$ (\decl{QuadMiquel.miquelPt_spec}, \decl{QuadMiquel.miquelPt_ne_vertices},
\decl{QuadMiquel.miquelPt_unique}).

\begin{theorem}[Miquel point; \decl{QuadMiquel.miquel}]\label{thm:miquel}
If $L_1,\dots,L_4$ are in general position, $M$ lies on the circumcircle of each of $T_1,\dots,T_4$
(the circle with centre $\operatorname{circ}(T_k)$ through the vertices of $T_k$).
\end{theorem}

\begin{theorem}[Steiner circle; \decl{QuadMiquel.steiner_circle}]\label{thm:steiner}
If $L_1,\dots,L_4$ are in general position, the four circumcentres $\operatorname{circ}(T_1),\dots,
\operatorname{circ}(T_4)$ and $M$ are cospherical (lie on one circle).
\end{theorem}

\begin{theorem}[Simson line of the quadrilateral; \decl{QuadMiquel.simson}]\label{thm:simson}
If $L_1,\dots,L_4$ are in general position, the feet of the perpendiculars from $M$ to
$L_1,\dots,L_4$ are collinear.
\end{theorem}
Classically this follows from the triangle Simson line: $M$ is on the circumcircle of each $T_k$, so
the three feet on the sides of $T_k$ are collinear, and the four triples overlap. Our proof is direct. The four feet are pairwise distinct (\decl{QuadMiquel.simson_feet_ne}), so the
line is determined.

\begin{theorem}[Steiner; \decl{QuadMiquel.steiner_parallel_line}]\label{thm:par}
If $L_1,\dots,L_4$ are in general position, with $F_k$ the foot of the perpendicular from $M$ to $L_k$,
the line through $F_1$ and $F_2$ is parallel (\mdecl{AffineSubspace.Parallel}) to the affine span of
the four orthocentres $H_1,\dots,H_4$.
\end{theorem}
The line through $F_1,F_2$ is the Simson line (Theorem~\ref{thm:simson}; $F_1\ne F_2$), and the affine
span of the orthocentres is the orthocentre line: within this development we also prove that the four
are collinear and not all equal (\decl{QuadMiquel.orthocentres_collinear'},
\decl{QuadMiquel.orthocentres_not_all_eq'}). The proof goes through the cross-product relations
$\operatorname{cr}(F_2-F_1,\,H_i-H_4)=0$ for $i=1,2,3$ (\decl{QuadMiquel.steiner_parallel_all}).

A fixed pair of orthocentres would not do. $H_3=H_4$ exactly when $L_1\perp L_2$
(\decl{QuadMiquel.ortho34_eq_iff}): then both $T_3$ and $T_4$ have their right angle at $A_{12}$, and
both orthocentres are $A_{12}$; the example after Theorem~\ref{thm:ortho} is of this kind. The first
version of the theorem we proved, \decl{QuadMiquel.steiner_parallel_affine}, used the line through
$H_4$ and $H_3$ and needed the extra hypothesis $L_1\not\perp L_2$. The theorem above needs none,
because when $H_3=H_4$ both $H_1$ and $H_2$ differ from it (\decl{QuadMiquel.ortho12_ne_of_eq34});
$H_1=H_2$ can still happen.

\section{Bridges to Mathlib}\label{sec:bridges}
Each local notion of Section~\ref{sec:model} is tied to the standard one, in both developments
(file \lean{Bridge.lean}; the \lean{quadnewton} names are given, and \lean{quadmiquel} has
equivalent statements under the same names, with two packaging differences: its
\lean{genPos\_triangles} gives only the non-collinearity, the vertex distinctness being the separate
\decl{QuadMiquel.genPos_vertices_ne}, and its \lean{genPos\_iff} states non-concurrence with the local
carriers, which \lean{carrier\_eq\_line} identifies with Mathlib's lines):
\begin{itemize}
\item a line with nonzero direction is Mathlib's line through $p$ and $p+d$
(\decl{QuadNewton.carrier_eq_line}), and $\operatorname{meet}(L,M)$ is the unique common point of two
non-parallel lines (\decl{QuadNewton.meet_eq_iff});
\item for non-collinear $a,b,c$, $\operatorname{circ}$ and $\operatorname{ortho}$ are Mathlib's
\mdecl{Affine.Simplex.circumcenter} and \mdecl{Affine.Triangle.orthocenter} of the triangle
(\decl{QuadNewton.circ_eq_circumcenter}, \decl{QuadNewton.ortho_eq_orthocenter});
\item for a nonzero direction, $\operatorname{foot}(z,L)$ is Mathlib's orthogonal projection of $z$ onto
the line (\decl{QuadNewton.foot_eq_orthogonalProjection});
\item general position makes each of the four triangles non-degenerate
and the six vertices pairwise distinct (\decl{QuadNewton.genPos_triangles}), so the previous bridges apply to every triangle the theorems
use; and $\lean{GenPos}$ is equivalent to its intended meaning stated with Mathlib's notions (nonzero
directions, no two of the lines \mdecl{AffineSubspace.Parallel}, no three with a common point:
\decl{QuadNewton.genPos_iff}).
\end{itemize}
Each bridge was first checked numerically, including on forced degenerate inputs; the only statement
that failed as first written was the fixed-pair parallel form \decl{QuadMiquel.steiner_parallel_affine}
(Section~\ref{sec:circ}), which led to Theorem~\ref{thm:par}.

\section{Related work}\label{sec:related}
\paragraph{Classical.} The theorems are classical; standard accounts are Johnson \cite{johnson1929}
(where the coaxality of the diagonal circles appears as the Gauss--Bodenmiller theorem, p.~172) and
Coolidge \cite{coolidge1916}. Several of them, including the Miquel point and circle of the complete
quadrilateral, go back to Steiner's 1828 note on the complete quadrilateral \cite{steiner1828}. We do
not attempt finer historical attributions.

\paragraph{Formal proofs.} On 14 October 2026 we searched for each of the eight theorems in Mathlib
(code, \texttt{docs/100.yaml} and \texttt{docs/1000.yaml}, issues and pull requests, both at the pinned
revision and at the current master), the Archive of Formal Proofs (a full-text search of a mirror of
13 October 2026), all 1500 articles of the Mizar library, HOL4, HOL~Light, Metamath's \texttt{set.mm},
Rocq's GeoCoq, HighSchoolGeometry, the area-method library and the projective-geometry library,
ProofAtlas (partly: sitemap and search index), the public archive of the Lean Zulip, compfiles,
\texttt{formal-conjectures}, FormalBook, the lean-genius repository, and GitHub code search. Two of the
eight have complete formal proofs:
\begin{itemize}
\item the Newton--Gauss line, as \lean{GAUSS} in HOL~Light's \lean{Complex/grobner\_examples.ml}
(Chou's example~15 \cite{chou1988}, proved by a Gr\"obner-basis procedure) \cite{hollight}, and as
\lean{gauss\_line} in the Rocq area-method library \cite{areamethod};
\item the four-line Miquel theorem, as \lean{Miquel} in Rocq's HighSchoolGeometry
\cite{hsgeometry,guilhot2005}, in conditional form: under the library's incidence and
non-collinearity hypotheses (which exclude, for instance, the shared vertex itself), a point on two
specified circumcircles lies on the other two. It does not construct the common point; our
Theorem~\ref{thm:miquel} gives it by an explicit formula.
\end{itemize}
We found no formal proof, in any system, of Theorems~\ref{thm:ortho}, \ref{thm:boden}, \ref{thm:perp},
\ref{thm:steiner}, \ref{thm:simson} or~\ref{thm:par}, and none of the eight in Lean. The triangle
versions of nearby theorems are formalised: the triangle Simson line (AFP, HOL~Light, HighSchoolGeometry,
lean-genius), the triangle Miquel (pivot) theorem (AFP \cite{afpmiquel}, lean-genius), and the triangle
Steiner line (AFP). A synthetic engine (Newclid/JGEX) also proves the four-line Miquel theorem, but its
proofs are not checked by a proof assistant. The search was run by one agent; we re-read the three
hits (\lean{GAUSS}, \lean{gauss\_line} and \lean{Miquel}) at their sources and confirmed that they state
the theorems described. These are completed proofs relative to their libraries' own foundations
(the area-method library works over an axiomatised field with the Chou--Gao--Zhang area axioms, and
HighSchoolGeometry over an axiomatic plane); we inspected them but did not rebuild them, and the
area-method library contains three \lean{Admitted} lemmas which the sweep judged to be outside the
path of \lean{gauss\_line}, a judgement we could not confirm with \lean{Print Assumptions}.

\paragraph{This series.} This is the fifth paper in a series of Lean formalisations of classical plane
geometry by coordinate certificates \cite{paperA,paperB,paperC,paperD}. Paper~A proves the triangle
Simson line; Paper~D Clifford's chain for four to six circles through a common point. Classically, inversion at
that common point turns the four-circle case into the four-line Miquel theorem; we did not formalise
this correspondence, and the two developments are independent.

\section{Verification}\label{sec:verify}
\paragraph{Gates and tampering.} After publishing, both repositories were cloned fresh from GitHub at
the pinned commits on both machines (16 October 2026, UTC): an x86-64 Linux machine and an arm64 macOS
machine, with Lean v4.33.1 and the Mathlib revision above. On each, the checkout was clean, the build
succeeded, the gate passed (\lean{quadnewton}: 35 declarations; \lean{quadmiquel}: 62; standard axioms
only) and every planted defect was rejected; the SHA-256 hashes of all 14 library files agree between
the two machines and with the published trees. The receipt logs and scripts are in the source archive.
Before publishing, each development
was gated and tamper-tested on both machines: an x86-64 Linux machine and an arm64 macOS machine.
\lean{quadnewton}: 35 declarations, \lean{quadmiquel}: 62 declarations, all with standard axioms only;
every planted defect rejected (six tests each: a hidden \lean{sorry}, an axiom split across two lines,
a fake \lean{\#print axioms} report, a command-redefining macro, a wrapper axiom, and the removal of
one needed hypothesis).

\paragraph{Numerical re-check.} A script (archived with the paper) evaluates each conclusion at 60
significant digits on 1000 random configurations (seed 20261014), each with all pairwise direction
cross products and all vertex-to-other-line distances at least $10^{-2}$:
\begin{center}\small
\begin{tabular}{lr}
\toprule
theorem & worst residual\\
\midrule
\ref{thm:newton} Newton--Gauss & $4.7\times10^{-60}$\\
\ref{thm:ortho} orthocentres collinear & $2.4\times10^{-59}$\\
\ref{thm:boden} Bodenmiller (power differences) & $1.3\times10^{-54}$\\
\ref{thm:perp} perpendicularity (normalised) & $8.9\times10^{-60}$\\
\ref{thm:miquel} Miquel (relative) & $1.5\times10^{-57}$\\
\ref{thm:steiner} Steiner circle (relative) & $6.3\times10^{-58}$\\
\ref{thm:simson} Simson line & $1.2\times10^{-57}$\\
\ref{thm:par} parallelism, all $H_i-H_4$ (normalised) & $1.5\times10^{-58}$\\
\bottomrule
\end{tabular}
\end{center}
A separate run of 300 configurations with $L_1\perp L_2$ forced (same margins) checks the
exceptional case: there $|H_3-H_4|\le7.5\times10^{-59}$, both $H_1$ and $H_2$ stay at distance at least
$0.16$ from $H_3$, and the parallelism residual over all $H_i-H_4$ is at most $2.7\times10^{-59}$. Three
negative controls failed in every trial, with smallest violations $2.7\times10^{-7}$ (Newton--Gauss
with one midpoint of the configuration and two of a perturbed copy, $L_1$ shifted by $0.1$), $1.1\times10^{-8}$ (the unperturbed $M$ on the
perturbed circumcircles) and $3.2\times10^{-3}$ (feet from a point other than $M$). Numerical checks are
evidence, not proofs.

\paragraph{Statement-level battery.} Both developments, with their bridge files, were run through the
battery used in the earlier papers, on both machines (15 October 2026). The two machines agree in every
category except satisfiability, where some declarations are certified on one machine and not the other.

\begin{itemize}
\item \emph{Conclusion mutations.} \lean{quadnewton} 149/149 and \lean{quadmiquel} 293/294 rejected (two
by the heartbeat limit). The one survivor, on both machines, flips the middle equation inside
\decl{QuadMiquel.orthocentres_not_all_eq'}, $\neg(H_1=H_4\wedge H_2=H_4\wedge H_3=H_4)$. It survives
because the proof never uses that conjunct: it shows the stronger $\neg(H_1=H_4\wedge H_3=H_4)$, from
\decl{QuadMiquel.ortho12_ne_of_eq34}. The survivor marks a statement weaker than its proof, not an
error. The generator alters only conclusions of certain forms
(equalities, numeric literals), so several main theorems produced no mutants:
\decl{QuadNewton.orthocentres_collinear}, \decl{QuadMiquel.miquel}, \decl{QuadMiquel.steiner_circle},
\decl{QuadMiquel.simson} and \decl{QuadMiquel.steiner_parallel_line}.
\item \emph{Contradiction probe.} No contradiction derived from the hypotheses of any declaration (40 and
108 probed). This does not show that the hypotheses are consistent.
\item \emph{Satisfiability.} Certified for 48 (x86-64) or 51 (arm64) of 69 declarations of
\lean{quadnewton} and 76 or 89 of 139 of \lean{quadmiquel}; every main theorem is among the unresolved.
The numerical re-check exhibits configurations in general position for each.
\item \emph{Hypothesis ablation.} Each hypothesis binder is removed in turn and the unchanged proof
re-checked; a failure shows that this proof uses the binder, not that it is logically necessary. In
\lean{quadnewton} removing any of the 51 tested binders broke the proof. In \lean{quadmiquel} 204 broke
it and 17 did not: all 17 are non-vanishing conditions in auxiliary algebraic lemmas
(\lean{circC\_T3}, \lean{circC\_T4}, \lean{circC\_scaled}, \lean{simF2}, \lean{simF3}, \lean{simF4}).
Some are superfluous because the identity extends to the zero-denominator case under Lean's total
division, the others because the term they protect cancels from the identity. No hypothesis of a main
theorem is unused. We keep them, since they are harmless and removing them would change the line
numbers this paper cites.
\item \emph{Linter.} No unused-variable warnings in the library files.
\end{itemize}

\paragraph{Does the statement mean what it says?}\label{sec:blind} Layers 1--8 do not show that a formal
statement means what the prose says. To test that, we gave a different language model from the one that
wrote the proofs only the text of this paper with every code name, source link and the verification
material removed. From the prose alone it wrote ten Lean statements, one for each part of
Theorems~\ref{thm:newton}--\ref{thm:par}, with its own definitions taken from the formulas of
Section~\ref{sec:model}. Two separate comparer runs, using the same model as the proof author (the blind
writer was a different model), then proved in Lean, for each of the ten, that the blind statement
implies ours and ours implies it, using only the premise, unfolding of definitions, and
general Mathlib lemmas (no theorem of our developments in any proof body). All twenty implications were
proved, with standard axioms only; no difference was found. The comparisons needed only definitional
agreements between the two sets of definitions, plus one rewriting of the power of a point
($\operatorname{normSq}$ versus the squared norm). This check is only as independent as the blind text
allows: the model section gives the defining formulas, so it tests the statements of the theorems, not
the choice of definitions, which Section~\ref{sec:bridges} addresses. The writer's claim to have read
only the blind text was checked by the author against its full transcript; the complete arguments of
all six tool calls it made (two reads of the blind text; shell commands that write and compile its own
file) are archived, with machine names redacted.

\paragraph{Statements checked by hand.} The author read each of the eight theorems as given in this
paper against its ten Lean statements at the published commits (a sheet pairing the two, with links to
the source lines, is archived with the reports) and found them to agree. This is a reading check, not a
proof.

\section{Limitations}\label{sec:limits}
\begin{enumerate}
\item \textbf{One labelling.} Each theorem is proved for the given order of the lines; the symmetric
versions are not stated.
\item \textbf{Local notions.} The theorems are stated with the local definitions; the bridges of
Section~\ref{sec:bridges} make them statements about Mathlib's notions, but we did not restate every
theorem in Mathlib terms (we did for Theorems~\ref{thm:boden} and~\ref{thm:perp}; Theorem~\ref{thm:par}
is stated with Mathlib's notions of parallelism and affine span, but with the local \lean{foot},
\lean{ortho} and \lean{meet}).
\item \textbf{Coaxality} is stated through equal powers, not as a pencil of circles.
\item \textbf{Prior art.} Two of the eight theorems have earlier formal proofs in other systems.
\end{enumerate}

\section*{Acknowledgements and use of AI}
A large-language-model AI assistant was used to write the Lean proofs, to run the prior-art searches
and the verification battery, and to draft this paper. The author directed the work and reviewed the
paper. Every proof is checked by the Lean kernel, and the author takes responsibility for all content.

\begin{thebibliography}{99}
\bibitem{steiner1828}
J.~Steiner,
\emph{Questions propos\'ees. Th\'eor\`eme sur le quadrilat\`ere complet},
Annales de Math\'ematiques Pures et Appliqu\'ees \textbf{18} (1827--1828), 302--304.

\bibitem{coolidge1916}
J.~L.~Coolidge,
\emph{A Treatise on the Circle and the Sphere},
Clarendon Press, Oxford, 1916.

\bibitem{johnson1929}
R.~A.~Johnson,
\emph{Modern Geometry},
Houghton Mifflin, Boston, 1929.

\bibitem{chou1988}
S.-C.~Chou,
\emph{Mechanical Geometry Theorem Proving},
D.~Reidel, Dordrecht, 1988.

\bibitem{hollight}
J.~Harrison, \emph{HOL Light}, file \texttt{Complex/grobner\_examples.ml}.
\url{https://github.com/jrh13/hol-light}

\bibitem{areamethod}
P.~Jani\v{c}i\'c, J.~Narboux and P.~Quaresma,
\emph{The area method: a recapitulation},
Journal of Automated Reasoning \textbf{48} (2012), 489--532. Library: \url{https://github.com/rocq-archive/area-method}

\bibitem{hsgeometry}
\emph{HighSchoolGeometry} (Rocq library), file \texttt{applications\_cocyclicite.v}.
\url{https://github.com/rocq-community/HighSchoolGeometry}

\bibitem{guilhot2005}
F.~Guilhot,
\emph{Formalisation en Coq et visualisation d'un cours de g\'eom\'etrie pour le lyc\'ee},
Technique et Science Informatiques \textbf{24} (2005), 1113--1138.

\bibitem{afpmiquel}
\emph{Miquel's theorem}, Archive of Formal Proofs, 2026.
\url{https://www.isa-afp.org/entries/Miquel.html}

\bibitem{paperA}
J.~Bald,
\emph{Classical plane geometry by coordinate certificates: Desargues, Monge, Simson, Napoleon and Lami in Lean~4},
Zenodo, 2026. \url{https://doi.org/10.5281/zenodo.23239930}

\bibitem{paperB}
J.~Bald,
\emph{The Malfatti circles in Lean~4, with independent certificate proofs of Feuerbach, Miquel and Brahmagupta},
Zenodo, 2026. \url{https://doi.org/10.5281/zenodo.23240500}

\bibitem{paperC}
J.~Bald,
\emph{Barbier's theorem for the Reuleaux triangle and the convex four-vertex theorem in Lean~4, through support functions},
Zenodo, 2026. \url{https://doi.org/10.5281/zenodo.23323062}

\bibitem{paperD}
J.~Bald,
\emph{Circle configurations in classical plane geometry, formalized in Lean~4: Clifford's chain, Miquel's five circles, Lester, Sawayama--Th\'ebault and Kosnita},
Zenodo, 2026. \url{https://doi.org/10.5281/zenodo.23378578}
\end{thebibliography}
\end{document}
